\begin{onehalfspace}

Il Capitolo~\ref{chap:sviluppo} ha descritto le scelte implementative dei due metodi: il Metodo 1 con strategia TOP-1 e il Metodo 2 con classificatore \ac{ML} e ricerca TOP-K. Questo capitolo mette alla prova l'ipotesi da cui è partita la parte sperimentale --- che un classificatore \ac{ML} applicato all'output rumoroso riduca il numero di iterazioni --- e ne documenta l'esito: \emph{che cosa è stato testato, perché, e che cosa è stato appurato}. Il risultato, anticipato qui per chiarezza, è netto: la riduzione delle iterazioni si ottiene, ma il merito non è del classificatore, bensì della semplice ricerca multi-candidato TOP-4 che lo accompagna; l'\ac{ML} applicato \emph{in quel punto e in quel modo} non serve. È questo esito, verificato con un'analisi di ablazione, a motivare il cambio di prospettiva che costituisce il nucleo della tesi --- la correzione degli errori \emph{durante} l'esecuzione --- presentato nei Capitoli~\ref{chap:qec}--\ref{chap:integrazione}.

% -----------------------------------------------
\section{Setup Sperimentale}
\label{sec:setup_m1}
\label{sec:setup_m2}

Tutti gli esperimenti sono stati condotti nell'ambiente WSL descritto nel Capitolo~\ref{chap:sviluppo}, con il simulatore \texttt{AerSimulator} in modalità \texttt{statevector} (CPU): \texttt{optimization\_level=2}, $M_{\text{shots}}=1024$ per iterazione, $M_{\text{max}}=50$ iterazioni, $K=30$ ripetizioni per UC1/UC2, seed deterministici.

\paragraph{Indipendenza fra le iterazioni.}
Un dettaglio dell'assegnazione dei seed merita di essere esplicitato, perché condiziona il significato stesso della metrica $\bar{M}$. Il simulatore deriva il seme di ciascuno shot dalla somma fra il seme dell'esecuzione e l'indice dello shot: due esecuzioni con semi consecutivi condividono quindi $1\,023$ campioni su $1\,024$ e costituiscono, in pratica, la medesima esecuzione. Poiché $\bar{M}$ intende misurare quante esecuzioni \emph{indipendenti} siano necessarie, le iterazioni devono essere separate di più del numero di shot; lo schema adottato è $\texttt{seed} = \text{rep} \times 10^6 + \text{iter} \times 10^4$. La verifica empirica conferma la necessità dell'accorgimento: la sovrapposizione fra gli istogrammi di due esecuzioni con semi consecutivi è del $99.9\%$, contro il $90.4\%$ fra semi ben separati --- valore, quest'ultimo, atteso fra campioni realmente indipendenti tratti da una distribuzione concentrata. I quattro use case seguono la struttura definita nel Capitolo~\ref{chap:obiettivi} (Tabella~\ref{tab:use_case_params}): UC1 e UC2 usano lo stesso circuito per $N=15$ a due livelli di rumore diversi (asse rumore), mentre UC3 ($N=21$) e UC4 ($N=35$) testano la scalabilità a rumore NISQ-realistico costante.

Le tre configurazioni algoritmiche confrontate sono:
\begin{itemize}
    \item \textbf{M1 (TOP-1)}: a ogni iterazione si tenta l'estrazione dei fattori dalla sola misura più frequente dell'istogramma. È la baseline, priva di parametri da ottimizzare.
    \item \textbf{M2 (CLF+TOP-4)}: un classificatore \ac{ML} pre-addestrato decide se l'istogramma contiene segnale \ac{QPE} recuperabile; in caso affermativo si tentano i 4 candidati più frequenti.
    \item \textbf{$M_{\text{TOP4}}$ (TOP-4 senza classificatore)}: variante di ablazione che applica la sola ricerca multi-candidato, senza filtro \ac{ML}. È il termine di confronto che isola il contributo del classificatore da quello della strategia di ricerca.
\end{itemize}

% -----------------------------------------------
\section{Analisi del Circuito e Barriera di Scalabilità}
\label{sec:circuito_m1}
\label{sec:scalabilita}

La profondità del circuito determina l'effetto cumulativo del rumore e definisce il perimetro entro cui il confronto è possibile.

\begin{table}[H]
\centering
\caption[Caratteristiche del circuito traspilato per ciascun use case]{Caratteristiche del circuito traspilato e applicabilità dei metodi per ciascun use case. I valori si riferiscono all'implementazione adottata (\texttt{c\_amod15} per $N=15$, Beauregard per $N=21$ e $N=35$), con \texttt{optimization\_level=2} e Qiskit 2.x. Fra parentesi, per confronto, il conteggio della sintesi densa \texttt{UnitaryGate} usata nella prima versione del codice.}
\label{tab:circuit_depth}
\resizebox{\textwidth}{!}{%
\begin{tabular}{lcccccc}
\hline
\textbf{Use Case} & \textbf{N} & \textbf{Qubit} & \textbf{Profondità} & \textbf{Porte CX} & \textbf{$P_{\text{surv}} = (1-\epsilon_{2q})^{\#\text{CX}}$} & \textbf{Metodi applicabili} \\
\hline
UC1 & 15 & 12 & 311    & 114    & $\approx 0.318$            & M1, $M_{\text{TOP4}}$, M2 \\
UC2 & 15 & 12 & 311    & 114    & $\approx 2.9\times10^{-3}$ & M1, $M_{\text{TOP4}}$, M2 \\
UC3 & 21 & 22 & 9\,883 & 3\,406 \ ($10\,040$) & $\approx 1.4\times10^{-15}$ & Nessuno \\
UC4 & 35 & 26 & 1\,480 & \phantom{0}516 \ ($13\,064$) & $\approx 5.6\times10^{-3}$ & Nessuno \\
\hline
\end{tabular}%
}
\end{table}

Per UC1 e UC2 il circuito traspilato conta 114 porte \texttt{cx} (profondità 311); la probabilità di sopravvivenza coerente $P_{\text{surv}}$ --- la frazione di run in cui nessun gate a 2 qubit introduce errore --- vale $\approx 31.8\%$ per UC1 e $\approx 0.29\%$ per UC2. Per UC3 e UC4 la tabella richiede due precisazioni, perché il motivo per cui nessun metodo risulta applicabile è diverso nei due casi e in nessuno dei due coincide con quello che una prima lettura suggerirebbe.

\textbf{UC3 è il caso genuinamente difficile.} Con $a=2$ e $N=21$ il periodo è $r=6$ e nessuno dei dieci gate controllati collassa nell'identità: il circuito è profondo per ragioni strutturali. La decomposizione di Beauregard~\cite{beauregard2002}, implementata nel modulo \texttt{beauregard.py}, porta il conteggio da $10\,040$ a $3\,406$ porte \ac{CX}, ma $P_{\text{surv}}$ resta dell'ordine di $10^{-15}$ e --- come mostra il paragrafo seguente --- il costo di simulazione diventa proibitivo proprio nel regime in cui la barriera fisica si attenuerebbe.

\textbf{UC4 è un caso di natura diversa}, e la tabella lo rende evidente in modo che merita di essere esplicitato. Con $a=6$ e $N=35$ il periodo è $r=2$, ossia $6^2 \equiv 1 \pmod{35}$: tutti i gate controllati $U^{2^j}$ con $j \geq 1$ sono l'identità e vengono eliminati dal transpiler, lasciando \emph{un solo} gate non banale su dodici. Il circuito risultante conta 516 porte \ac{CX}, meno di un sesto di UC3 pur avendo un modulo più grande, e $P_{\text{surv}} \approx 5.6\times10^{-3}$ --- un valore che la campagna parametrica su UC1 colloca ampiamente nel regime in cui la strategia multi-candidato funziona. UC4 non è dunque fuori portata per il rumore: lo è soltanto per il costo di simulazione dei suoi 26 qubit, che rende ogni traiettoria proibitiva.

Questo esito è coerente con il razionale dichiarato nel Capitolo~\ref{chap:obiettivi}, dove $r=2$ era stato scelto \emph{deliberatamente} per separare l'effetto della larghezza del registro da quello della profondità dell'esponenziazione modulare. La scelta è legittima e il risultato la conferma: UC4 misura la prima variabile a profondità minima. Ma proprio per questo UC4 \textbf{non appartiene all'asse di scalabilità} lungo cui l'Ipotesi secondaria 2 lo colloca. La difficoltà del circuito non cresce con $N$: cresce con il periodo $r$, che determina quanti dei gate controllati sopravvivono alla traspilazione. Lungo la sequenza $N = 15 \to 21 \to 35$ i periodi sono $4 \to 6 \to 2$, e con essi la profondità va $311 \to 9\,883 \to 1\,480$: non è una progressione monotòna, e trattarla come tale è un difetto di formulazione dell'ipotesi, non del disegno degli use case.

Ne discende una conseguenza che va dichiarata subito, perché riguarda il perimetro di validità di tutto il capitolo: l'\textbf{Ipotesi secondaria 2} formulata nel Capitolo~\ref{chap:obiettivi} --- che il vantaggio del Metodo 2 si mantenga al crescere dell'istanza ($N = 15 \to 21 \to 35$) --- \textbf{resta non verificata}, e va inoltre riformulata. Non è stata smentita: su UC3, l'unico dei tre che aumenti davvero la difficoltà, la verifica è impedita da una barriera doppia, fisica e computazionale; su UC4 dal solo costo di simulazione, e comunque su un'istanza che per costruzione non mette alla prova la profondità. Il piano sperimentale prevedeva quattro use case come requisito di validazione; ne restano misurabili due. Il Capitolo~\ref{chap:sviluppi} riprende il punto indicando come andrebbe riformulata l'ipotesi perché l'asse di scalabilità sia quello effettivo.

\paragraph{La barriera è doppia, e la decomposizione ne rimuove solo metà.}
Poiché l'attribuzione della barriera alla sola decomposizione sarebbe incompleta, la si è verificata direttamente. Il modulo \texttt{beauregard.py} mantiene la promessa: per $N=21$ il circuito passa da $10\,040$ a $3\,406$ porte \ac{CX} su 22 qubit (profondità $9\,883$), una riduzione di un fattore $2.9$. Ciò nondimeno UC3 resta irraggiungibile, per due ragioni indipendenti che è utile separare.

La prima è \textbf{fisica}: al livello di rumore previsto per UC3 ($\varepsilon_{2q} = 10^{-2}$) la probabilità di sopravvivenza coerente vale $P_{\text{surv}} = (1-10^{-2})^{3406} \approx 1.4\times10^{-15}$. Anche disponendo di una simulazione istantanea, l'output resterebbe presumibilmente indistinguibile dal rumore uniforme: la decomposizione efficiente non basta, occorre \emph{anche} un errore di gate più basso. A $\varepsilon_{2q} = 10^{-3}$ la stessa espressione dà $P_{\text{surv}} \approx 3.3\%$, valore con cui la strategia multi-candidato è plausibilmente applicabile.\footnote{L'avverbio è d'obbligo: la campagna parametrica su UC1 (Sezione~\ref{sec:parametri_mitigazione}) ha mostrato che $P_{\text{surv}}$ \emph{sottostima} sistematicamente la recuperabilità del segnale --- a $P_{\text{surv}} = 6\times10^{-6}$ la strategia TOP-4 converge ancora al primo tentativo. La stima analitica indica quindi un ordine di grandezza, non una soglia di impossibilità; solo la seconda barriera, quella computazionale, impedisce di verificarlo per UC3.}

La seconda è \textbf{computazionale}, e si manifesta proprio nel regime in cui la prima si attenua. In presenza di un modello di rumore il simulatore non può evolvere lo stato una volta sola e campionarlo: ogni shot richiede una traiettoria indipendente. Le misure condotte su questo circuito danno $\approx 100\,$s per una singola evoluzione priva di rumore e $\approx 50\,$s \emph{per shot} in presenza di rumore, da cui circa quattordici ore per una sola iterazione da $1\,024$ shot --- mentre il protocollo ne richiede fino a cinquanta per ciascuna delle trenta ripetizioni. Il metodo \textit{matrix product state}, che su circuiti poco entangled sarebbe più efficiente, risulta qui ancora più lento: l'algoritmo di Shor massimizza l'entanglement per costruzione, e la dimensione di legame cresce fino a rendere il vantaggio nullo.

La formulazione corretta dell'esito è dunque la seguente: \emph{l'Ipotesi secondaria 2 richiede un errore di gate a due qubit di circa $10^{-3}$ per essere anche solo formulabile, e a quel livello la verifica per simulazione classica diventa proibitiva}. Non è un limite degli algoritmi di analisi né della decomposizione, ma la conseguenza congiunta della fisica del dispositivo e del costo della simulazione: la sua verifica appartiene propriamente all'hardware reale.

\begin{figure}[H]
\centering
\includegraphics[width=0.65\textwidth]{12_QPE_circuit_N21.pdf}
\caption[Circuito QPE per $N=21$]{Circuito di Quantum Phase Estimation per $N=21$, $a=2$ ($n_{\text{count}}=10$ qubit di controllo, 5 qubit target). La complessità della decomposizione \texttt{UnitaryGate} porta a 10\,040 porte CX dopo la traspilazione.}
\label{fig:qpe_circuit_N21}
\end{figure}

\begin{figure}[H]
\centering
\includegraphics[width=0.7\textwidth]{13_histogram_simulator_N21.pdf}
\caption[Istogramma QPE per $N=21$ — distribuzione piatta]{Distribuzione di misura del registro di controllo per $N=21$ con decomposizione \texttt{UnitaryGate} (1024 shot, noise model NISQ-realistico): l'output è indistinguibile da una distribuzione uniforme.}
\label{fig:histogram_N21_flat}
\end{figure}

% -----------------------------------------------
\section{La Baseline: Metodo 1 (TOP-1)}
\label{sec:risultati_uc_m1}
\label{sec:analisi_m1}
\label{sec:uc1_m1}
\label{sec:uc2_m1}

La Figura~\ref{fig:istogrammi_qpe} mostra i due volti dell'output di UC1: un'iterazione con segnale QPE strutturato (picchi intorno ai valori teorici $\{64, 128, 192\}$) e una in cui il rumore rende l'output quasi uniforme. Distinguere i due casi \emph{prima} della post-elaborazione è esattamente il compito che l'ipotesi iniziale affidava al classificatore.

\begin{figure}[H]
    \centering
    \includegraphics[width=0.9\textwidth]{fig_istogrammi_qpe.pdf}
    \caption[Confronto istogrammi QPE con segnale strutturato e dominato dal rumore (UC1)]{Confronto tra un istogramma QPE con segnale strutturato (sinistra) e uno dominato dal rumore (destra) per UC1 ($N=15$, NISQ-realistico). I valori teorici attesi per $r=4$ sono $\{64, 128, 192\}$ su $2^8 = 256$ stati.}
    \label{fig:istogrammi_qpe}
\end{figure}

\begin{table}[H]
\centering
\caption{Risultati del Metodo 1 sui quattro use case}
\label{tab:riepilogo_m1}
\begin{tabular}{lcccc}
\hline
\textbf{Use Case} & \textbf{$P_{\text{succ}}$} & \textbf{$\bar{M}_1$} & \textbf{$\sigma_{M_1}$} & \textbf{Mediana} \\
\hline
UC1 ($N=15$, realistico) & 100\% (30/30) & 1.97 & 1.66 & 1.5 \\
UC2 ($N=15$, degradato)  & 100\% (30/30) & 1.60 & 0.76 & 1.0 \\
UC3 ($N=21$)             &  0\%   & N/A  &  ---  & --- \\
UC4 ($N=35$)             &  0\%   & N/A  &  ---  & --- \\
\hline
\end{tabular}
\end{table}

La distribuzione del numero di iterazioni è asimmetrica e approssimativamente geometrica: per UC1 ogni iterazione ha probabilità indipendente $p_{\text{iter}} = 58.3\%$ di produrre la misura corretta tramite TOP-1 (misura diretta su 60 esecuzioni), da cui $\bar{M}_1 \approx 2$. Il dato saliente è che \textbf{la baseline non fallisce mai}: entro il budget di cinquanta iterazioni tutte e trenta le ripetizioni convergono, su entrambi gli use case. La differenza fra i metodi non riguarda dunque \emph{se} il risultato venga trovato, ma \emph{quante esecuzioni} servano per trovarlo --- ed è precisamente ciò che la metrica $\bar{M}$ misura.

Il confronto UC1 vs UC2 rivela un comportamento controintuitivo ($\bar{M}_1$ \emph{minore} sotto rumore maggiore), spiegato dalla banda di accettazione di \texttt{extract\_factors} (\texttt{Fraction.limit\_denominator}): il rumore allarga la distribuzione e aumenta la probabilità che la moda ricada nella banda.

La ragione per cui TOP-1 richiede più di un tentativo è strutturale, e vale la pena isolarla perché anticipa il risultato della sezione seguente. Su sessanta esecuzioni indipendenti di UC1 la moda dell'istogramma cade su uno solo di quattro valori --- $\{0,\,64,\,128,\,192\}$, ossia esattamente i picchi $j\,2^{8}/r$ attesi dalla \ac{QPE} per $r=4$ --- e coincide con l'esito $y=0$ in $25$ casi su $60$ ($41.7\%$), che ne fa l'esito \emph{più frequente in assoluto}. Ma $y=0$ corrisponde a fase nulla e non porta alcuna informazione sul periodo: la funzione di estrazione lo scarta per costruzione. In altri termini, il candidato che TOP-1 seleziona è, in due casi su cinque, quello strutturalmente privo di contenuto informativo; la moda conduce ai fattori nel $58.3\%$ delle esecuzioni, da cui $\bar{M}_1 \approx 1/0.583 = 1.7$, in accordo con il valore misurato $1.97$. Quando invece si considerano i quattro candidati più frequenti, i picchi utili sono presenti in \emph{tutte} e sessanta le esecuzioni: è questa la ragione per cui $M_{\text{TOP4}}$ converge al primo tentativo con deviazione standard nulla.

I limiti della baseline sono dunque due, e motivano il passo successivo: (i) \textbf{spreco informativo} --- TOP-1 usa una sola misura dell'istogramma e scarta tutto il resto, incluso il caso in cui la misura selezionata sia l'esito banale; (ii) \textbf{variabilità} --- $\sigma_{M_1} = 1.66$ contro $\bar{M}_1 = 1.97$: il costo di una sessione resta imprevedibile a priori e non si riduce aumentando gli shot, poiché dipende da quale picco risulti dominante in quella specifica realizzazione del rumore.

% -----------------------------------------------
\section{L'Ipotesi ML alla Prova}

L'ipotesi formalizzata nel Capitolo~\ref{chap:obiettivi} attribuiva la riduzione delle iterazioni a un classificatore \ac{ML}, addestrato a riconoscere gli istogrammi privi di segnale QPE ($\bar{M}_2 \approx 3$, $\rho = \bar{M}_1/\bar{M}_2 > 1$). Sono state addestrate e confrontate tre architetture (Random Forest, \ac{SVM}, \ac{MLP}) su un dataset di 2000 istogrammi per use case: come \emph{classificatori} i modelli funzionano molto bene (F1 fino a 0.98, \ac{AUC} fino a 0.99). L'architettura, la procedura di addestramento e le metriche complete --- con la tabella dei tre modelli e le curve ROC --- sono riportate nell'Appendice~\ref{app:classificatore} (Sez.~\ref{sec:valutazione_clf}). Le metriche di classificazione, però, dicono soltanto che il modello sa distinguere le due classi di istogrammi; la domanda operativa è un'altra --- \emph{questo filtraggio riduce davvero le iterazioni rispetto alla sola ricerca multi-candidato TOP-4?} --- ed è quella a cui risponde l'ablazione seguente.

% -----------------------------------------------
\section{Il Verdetto dell'Ablazione}
\label{sec:confronto}
\label{sec:uc1_m2}
\label{sec:uc2_m2}

\begin{table}[H]
\centering
\caption{Riepilogo comparativo M1, $M_{\text{TOP4}}$ e M2 --- risultato principale con analisi di ablazione ($K=30$, $M_{\text{max}}=50$)}
\label{tab:riepilogo_rho}
\begin{tabular}{llcccc}
\hline
\textbf{UC} & \textbf{Variante} & \textbf{$P_{\text{succ}}$} & \textbf{$\bar{M}$} & \textbf{$\rho$ vs M1} & \textbf{p-value (vs M1)} \\
\hline
\multirow{3}{*}{UC1} & M1 (TOP-1)                 & 100\% & 1.97 & ---            & --- \\
                     & $M_{\text{TOP4}}$ (no clf) & 100\% & 1.00 & \textbf{1.967} & ${<}0.001$ (sign.) \\
                     & M2 (CLF+TOP-4)             & 100\% & 1.03 & 1.903          & ${<}0.001$ (sign.) \\
\hline
\multirow{3}{*}{UC2} & M1 (TOP-1)                 & 100\% & 1.60 & ---            & --- \\
                     & $M_{\text{TOP4}}$ (no clf) & 100\% & 1.00 & \textbf{1.600} & ${<}0.001$ (sign.) \\
                     & M2 (CLF+TOP-4)             & 100\% & 1.57 & 1.021          & $0.315$ (n.s.) \\
\hline
\end{tabular}

\smallskip
{\footnotesize Ablazione $M_{\text{TOP4}}$ vs M2 (Mann-Whitney): UC1 $p = 0.849$ (n.s.); UC2 $p < 0.001$, con M2 che richiede significativamente \emph{più} iterazioni di $M_{\text{TOP4}}$.}
\end{table}

\paragraph{UC1: il classificatore è superfluo.}
M2 raggiunge $\bar{M}_2 = 1{,}03$ con $\rho = 1{,}903$ significativo rispetto a M1: l'ipotesi numerica ($\bar{M}_2 \approx 3$) è addirittura superata. Ma $M_{\text{TOP4}}$ --- TOP-4 \emph{senza} classificatore --- ottiene un risultato migliore ($\bar{M} = 1{,}00$ esatto, con deviazione standard nulla: \emph{ogni} ripetizione converge al primo tentativo), e il confronto $M_{\text{TOP4}}$ vs M2 non è significativo ($p = 0.849$). Il guadagno è interamente della ricerca multi-candidato: il picco QPE si trova sempre tra i quattro candidati più frequenti anche sotto rumore, perché il rumore allarga la distribuzione ma non ne distrugge la struttura --- e, come osservato sopra, il candidato che TOP-1 scarta a favore dell'esito banale $y=0$ viene qui recuperato.

\paragraph{UC2: il classificatore annulla il guadagno.}
M2 produce $\bar{M}_2 = 1{,}57$ contro $\bar{M} = 1{,}00$ del solo $M_{\text{TOP4}}$: il classificatore riporta il costo quasi al livello della baseline ($\rho = 1{,}021$, non significativo rispetto a M1), vanificando il vantaggio che la ricerca multi-candidato aveva ottenuto. Il classificatore \ac{SVM}, addestrato con rumore variato $\pm50\%$, risulta eccessivamente conservativo al rumore nominale di UC2 e scarta iterazioni con segnale QPE parziale che la sola ricerca TOP-4 avrebbe risolto. Rispetto a UC1, dove il filtro è semplicemente inerte, qui esso è attivamente controproducente rispetto alla strategia che accompagna.

\begin{figure}[H]
    \centering
    \includegraphics[width=0.75\textwidth]{fig_iterazioni_m1_m2.pdf}
    \caption[Analisi di ablazione: M1, $M_{\text{TOP4}}$ e M2 su UC1 ($K=30$)]{Analisi di ablazione su UC1 ($K=30$, $M_{\text{max}}=50$): confronto tra M1 (TOP-1), $M_{\text{TOP4}}$ (TOP-4 senza classificatore) e M2 (CLF+TOP-4). $M_{\text{TOP4}}$ è statisticamente indistinguibile da M2 ($p=0.849$): il guadagno è della strategia TOP-4, non del classificatore.}
    \label{fig:iterazioni_m1_m2}
\end{figure}

\paragraph{Che cosa è stato appurato.}
Il risultato qualitativo dell'ipotesi --- che un approccio multi-candidato riduca $\bar{M}$ --- è confermato per entrambi gli use case tramite $M_{\text{TOP4}}$. La parte dell'ipotesi che attribuiva il guadagno al classificatore è \textbf{smentita}: contributo nullo su UC1, negativo su UC2. Le eccellenti metriche di classificazione non si traducono in guadagno operativo perché \emph{l'informazione che il classificatore estrae dall'istogramma è la stessa che la ricerca TOP-4 sfrutta direttamente}, senza addestramento e senza rischio di falsi negativi. Applicare l'\ac{ML} lì --- a valle dell'esecuzione, sul solo output --- e in quel modo, non ha senso.

% -----------------------------------------------
\section{Robustezza della Strategia TOP-4}
\label{sec:parametri_mitigazione}
\label{subsec:sintesi_parametri}

Identificato il componente attivo, una campagna parametrica sistematica a variazione singola ($K_{\text{rep}}=30$ per punto, test Mann-Whitney) ne ha caratterizzato i margini operativi su sette parametri. Il dettaglio completo, con il confronto previsione/esito per ciascun parametro, è riportato nell'Appendice~\ref{app:parametrica}; la Tabella~\ref{tab:sintesi_parametri} ne riassume gli esiti.

\begin{table}[H]
\centering
\caption{Sintesi della campagna parametrica su $M_{\text{TOP4}}$: impatto atteso, osservato e raccomandazioni}
\label{tab:sintesi_parametri}
\resizebox{\textwidth}{!}{%
\begin{tabular}{lllll}
\hline
\textbf{Parametro} & \textbf{Categoria} & \textbf{Impatto atteso} & \textbf{Impatto osservato} & \textbf{Raccomandazione} \\
\hline
$K$                   & Strategia & Alto ($K = r = 4$)  & Alto ma soglia a $K=2$, non $K=4$              & $K = 2$ minimo; $K = 4$ per margine \\
$\varepsilon_{2q}$    & Rumore    & Alto (dominante)    & Nullo su $M_{\text{TOP4}}$; soglia non trovata  & TOP-K robusto su $[10^{-3}, 10^{-1}]$ \\
$M_{\text{shots}}$    & Strategia & Medio (soglia 1024) & Nullo; 128 shot già sufficienti                & Nessun requisito minimo pratico \\
$T_1$, $T_2$          & Rumore    & Basso (per $N=15$)  & Nullo su $[20\,\mu\text{s}, 500\,\mu\text{s}]$ & Irrilevante per $N=15$ \\
$\varepsilon_{1q}$    & Rumore    & Basso               & Nullo su $[10^{-4}, 2\times10^{-2}]$           & Trascurabile rispetto a $\varepsilon_{2q}$ \\
$p_{\text{ro}}$       & Rumore    & Basso               & Nullo su $[0\%, 20\%]$                         & Irrilevante su tutto il range testato \\
\hline
\end{tabular}%
}
\end{table}

Su sette parametri, uno solo --- $K$ --- ha impatto, con soglia effettiva $K=2$ (più bassa del valore teorico $K=r=4$); tutti gli altri, incluso $\varepsilon_{2q}$ su tre ordini di grandezza, risultano irrilevanti. $M_{\text{TOP4}}$ è quindi applicabile su hardware NISQ senza alcuna calibrazione: il meccanismo non si basa sulla preservazione del segnale QPE \emph{in media} ma sulla sua recuperabilità \emph{per realizzazione} --- come mostra il caso estremo $\varepsilon_{2q}=10^{-1}$, dove la probabilità che nessuna porta a due qubit sbagli è di sei parti su un milione e ciononostante ogni ripetizione converge al primo tentativo.

Il confronto con la \ac{ZNE} sullo stesso task (Tabella~\ref{tab:confronto_zne}; dettaglio in Appendice~\ref{app:parametrica}) completa il quadro. Nessuna delle due varianti di estrapolazione degrada il risultato --- tutte le strategie raggiungono il $100\%$ di successo --- ma nessuna lo migliora in modo significativo: $\bar{M}$ resta statisticamente indistinguibile da quello della baseline, mentre il costo in shot raddoppia o triplica per il campionamento ai livelli di rumore amplificato. La \ac{ZNE} è dunque \emph{inefficace}, non controproducente: paga un prezzo certo per un beneficio che su questo compito non si manifesta. La strategia che accetta il rumore e amplia la ricerca, al contrario, dimezza le iterazioni al costo di un solo campionamento.

\begin{table}[H]
\centering
\caption{Confronto M1, ZNE-2, ZNE-3 e $M_{\text{TOP4}}$ su UC1 ($K_{\text{rep}}=30$, $M_{\text{shots}}=1024$ per livello)}
\label{tab:confronto_zne}
\resizebox{\textwidth}{!}{%
\begin{tabular}{lcccccc}
\hline
\textbf{Strategia} & \textbf{$P_{\text{succ}}$} & \textbf{$\bar{M}$} & \textbf{Shot/iter} & \textbf{Shot totali} & \textbf{$\rho$ vs M1} & \textbf{p-value} \\
\hline
M1 (baseline)             & 100\% & 1.97 & 1024  & 2\,014 & ---   & --- \\
ZNE-2 (lineare)           & 100\% & 2.03 & 2048  & 4\,164 & 0.967 & $0.728$ (n.s.) \\
ZNE-3 (Richardson)        & 100\% & 1.57 & 3072  & 4\,813 & 1.255 & $0.250$ (n.s.) \\
$M_{\text{TOP4}}$ ($K=4$) & 100\% & 1.00 & 1024  & 1\,024 & 1.967 & ${<}0.001$ (sign.) \\
\hline
\end{tabular}%
}
\end{table}

% -----------------------------------------------
\section{Bilancio della Verifica e Ponte verso il Nucleo della Tesi}
\label{sec:top4_mitigazione}
\label{sec:validita}
\label{sec:confronto_zne}

Questa prima fase sperimentale ha appurato tre fatti, tutti con significatività statistica e piena riproducibilità (seed, noise model, circuito e split espliciti e fissati; codice nella repository della tesi):

\begin{enumerate}
    \item la ricerca multi-candidato TOP-4 riduce $\bar{M}$ da 1.97 a 1.00 (UC1), portando ogni ripetizione a convergere al primo tentativo, ed è robusta su tutto il range di rumore NISQ testato, senza addestramento e senza calibrazione;
    \item il classificatore \ac{ML} a valle dell'esecuzione non aggiunge nulla (UC1) o peggiora (UC2): l'ipotesi iniziale, nella parte che attribuiva il guadagno all'apprendimento automatico, era sbagliata --- e l'ablazione lo dimostra;
    \item le tecniche di correzione dell'output (ZNE) sono inefficaci su questo task: non degradano il risultato, ma non lo migliorano in modo significativo e costano da due a tre volte più shot.
\end{enumerate}

Il punto 2 merita di essere letto in positivo, perché delimita con precisione \emph{dove} si ferma ogni strategia a valle: sul solo output, l'informazione utile è già interamente catturata da una ricerca multi-candidato a costo zero --- e ciò che il rumore ha distrutto durante l'esecuzione non è più recuperabile da nessun post-processing, per quanto sofisticato. La conseguenza è strutturale: per andare oltre, l'errore va intercettato e corretto \emph{mentre} il calcolo è in corso, agendo sullo stato quantistico prima che l'informazione sia persa. È esattamente il programma dei prossimi capitoli, che costituiscono il nucleo di questa tesi: la correzione d'errore quantistica --- dal codice a ripetizione al codice di Steane (Capitolo~\ref{chap:qec}), al surface code con la stima sperimentale dell'errore logico (Capitolo~\ref{chap:surface}) --- fino allo Shor logico, che quantifica il livello di correzione necessario all'algoritmo (Capitolo~\ref{chap:integrazione}).

\end{onehalfspace}
